\begin{problem}{보안 등급}{standard input}{standard output}{1 second}{256 megabytes}

서버는 시간순으로 주어진 $N$개의 이벤트를 순서대로 검사한다.
각 이벤트 $i$에는 신뢰도 $A_i$가 주어진다.

서버는 하나의 정수 기준값 $T$를 정한 뒤, 각 이벤트를 다음과 같이 분류한다.

\begin{itemize}
    \item $A_i \ge T$이면 그 이벤트는 \textbf{안전 상태}가 된다.
    \item $A_i < T$이면 그 이벤트는 \textbf{위험 상태}가 된다.
\end{itemize}

서버에는 보안 \texttt{credit}이 있으며, 초기값은 $0$이다.
이벤트를 하나 검사할 때마다 보안 \texttt{credit}은 다음과 같이 변한다.

\begin{itemize}
    \item 이벤트가 안전 상태이면 \texttt{credit}이 $1$ 증가한다.
    \item 이벤트가 위험 상태이면 \texttt{credit}이 $1$ 감소한다.
\end{itemize}

검사 도중 어느 순간이라도 \texttt{credit}이 음수가 되면 서버는 즉시 다운된다.

서버는 총 예산 $B$를 가지고 있으며, 원래는 위험 상태인 이벤트를 패치하여 안전 상태로 바꿀 수 있다.

이벤트 $i$가 원래 $A_i < T$라면, 그 이벤트를 검사하는 순간에 패치할 수 있다.
이 경우 그 이벤트는 안전 상태로 간주되며, 필요한 비용은 $T - A_i$이다.

패치 비용의 총합은 $B$를 초과할 수 없다.

모든 이벤트를 끝까지 검사하는 동안 서버가 한 번도 다운되지 않도록 할 수 있는 최대의 정수 기준값 $T$를 구하여라.

\InputFile
첫째 줄에 이벤트의 개수 $N$과 예산 $B$가 주어진다.
$(1 \le N \le 200\,000,\ 0 \le B \le 10^9)$

둘째 줄에 $N$개의 정수 $A_1, A_2, \dots, A_N$이 주어진다.
$(1 \le A_i \le 10^9)$

\OutputFile
조건을 만족하도록 선택할 수 있는 최대의 정수 보안 등급 $T$를 출력한다.

\Example

\begin{example}
\exmpfile{example.01}{example.01.a}%
\end{example}

\end{problem}

